\section{Metodologia}
\subsection{Objetivo}
% TODO: comparação qualitativa e quantitativa ou apenas qualitativa
O presente estudo busca realizar uma comparação qualitativa entre diferentes
métodos para minerar padrões de uso de APIs da \textit{glibc} a fim de
proporcionar uma visão clara das vantagens e desvantagem de cada forma de
solucionar o problema. Busca-se também disponibilizar publicamente um catálogo
de padrões de uso no mesmo domínio para uso posteriormente em outros estudos. 

Em específico serão avaliados os métodos: 

\begin{itemize}
  \item aprendizado de regras de associação
  \item clustering
  \item modelo probabilístico
\end{itemize}

A seguir é especificado como tais são aplicados ao problema de mineração de
padrões de uso.

\subsubsection{Aprendizado de Regras de Associação}

Cada função transformada em uma transação e os itens nela contidos são as
funções que esta chama. O identificador único das funções é seu nome acrecido do
caminho até o arquivo que o contem e o nome do projeto no qual este arquivo se
encontra.

\subsubsection{Clustering}

\subsubsection{Modelo Probabilístico}

\subsection{Avaliação dos Métodos}

% TODO: é necessário apresentar como a lista de repositórios foi obtida ?
Todos os algoritmos apresentados nas seções anteriores são aplicados sobre o
mesmo conjunto de dados, uma lista de projetos disponibilizados no Github
contendo um número mínimo de contribuidores / avaliações afim de obter apenas
códigos que tenham alguma garantia de qualidade. Antes da execução dos
algoritmos, contudo, uma série de etapas de pré-processamento são realizadas,
como é demonstrado na figura \ref{fig:pipeline_overview}.

%TODO: traduzir imagem
\begin{figure}[h]
  \includegraphics[width=1\textwidth]{pipeline_overview}
  \caption{Visão global das etapas para minerar os padrões de uso neste estudo}
  \label{fig:pipeline_overview}
\end{figure}

A partir dos arquivos fontes escritos em na linguagem C nos projetos
selecionados, um grafo de chamadas é extraído para as análises. A fim de limitar
o escopo de tal dataset a apenas o escopo do presente estudo, apenas a chamadas
de funções da \textit{glibc} são levadas em conta. Assim, o grafo de chamadas
minerado dos arquivos é filtrado para se adequar a tal.

Porém, para realizar tal redução, faz-se necessário obter uma lista das funções
definidas na \textit{glibc}. Tal tarefa, como é visível na
figura \ref{fig:pipeline_overview} é feita em paralelo a tarefa de minerar o
grafo de chamadas.

Uma vez que estes dois processos estão finalizados, os dados podem são então
formatados para adequar-se a entrada esperada por cada um dos algoritmos
utilizados.
